Forecasts of glucose 60 minutes ahead from continuous glucose monitors (CGM) can warn people with type 1 diabetes before glucose becomes dangerously low or high. A forecast is more useful when it comes with a prediction interval, or band, that states how far it can be trusted. Bands are usually judged by their average coverage, the share of true values that fall inside them, and an average can hide failures in particular situations and for particular patients. This initial draft asks whether a nominal 80\% band stays honest in each situation, defined as the current glucose level combined with its 30-minute trend. Per-patient forecasters were trained for 215 people from 11 public studies, and eight of the studies (153 patients) were held out for external validation. Pooled over all studies, a standard GARCH band covered between 62\% and 91\% of outcomes across the situations with enough data. A situation-aware band, which learns a separate width for each situation, kept coverage between 72\% and 82\% and reduced per-situation miscalibration in all eight held-out studies while producing narrower bands. These results rest on one training seed per patient and are preliminary. Analyses by patient and by day are still in progress.
